% GENERATED FILE. Edit canonical lessons, then run npm run generate:books.
% canonical-source-hash: fnv1a64:abbc86b5ee06fda6
% canonical-lessons: BN-C494-bekar, BN-C494-bayodata, BN-C494-shikkhanobish, BN-C494-intarbhiu, BN-C494-jog-deoya

\chapter{Looking for Work}
\label{ch:bn-looking-for-work}
\hlchaptermodality{494}

\begin{chapteropening}
\textbf{By the end of this chapter:} \emph{I can talk about looking for work, CVs and joining a new job.}

Complete the last lesson of chapter 494: I can talk about looking for work, CVs and joining a new job.
\end{chapteropening}

\section[\texorpdfstring{\bn{বেকার} (bekār)}{বেকার (bekār)}]{\bn{বেকার} (bekār) — unemployed; out of work}

\label{lesson:BN-C494-bekar}



\begin{quote}
Before the new one: say the Bengali for an office messenger, then the Bengali for to type.
\end{quote}

\subsection*{You'll want to know: \bn{বেকার}}
\textbf{\bn{বেকার}} — \emph{bekār} — \textquotedblleft{}unemployed; out of work\textquotedblright{}.

\begin{quote}
\bn{আমার} \bn{ভাই} \bn{এক} \bn{বছর} \bn{ধরে} \bn{বেকার।}
\end{quote}

\begin{quote}
\emph{Āmār bhāi ek bôchhôr dhôre bekār.}
\end{quote}

\begin{quote}
My brother has been out of work for a year.
\end{quote}

\begin{cousinweb}
\textbf{\bn{বেকার}} (\emph{bekār}) is usually traced to Persian \emph{bekār}, \textquotedblleft{}without work\textquotedblright{}: \emph{be-}, \textquotedblleft{}without\textquotedblright{}, and \emph{kār}, \textquotedblleft{}work\textquotedblright{}.
\end{cousinweb}

\begin{grammarlens}[title={What you've built}]
One more word for describing what you lived and felt.
\end{grammarlens}

\subsection*{Your turn}
\begin{itemize}
\raggedright
  \item \emph{Say it:} \emph{bekār}
  \item \emph{Say it:} \emph{bekār}, once more
  \item \emph{Say it:} \emph{Āmār bhāi ek bôchhôr dhôre bekār.} — My brother has been out of work for a year.
  \item \emph{Recall:} say the Bengali for an office messenger, then the Bengali for to type, then say \emph{bekār} again
\end{itemize}

\begin{tcolorbox}[breakable,colback=teal!4,colframe=teal!35!black,title={Before you move on}]
What does \bn{বেকার} mean? (\textquotedblleft{}Unemployed; out of work\textquotedblright{}.) Say it once more.
\end{tcolorbox}
\section[\texorpdfstring{\bn{বায়োডাটা} (bāyoḍāṭā)}{বায়োডাটা (bāyoḍāṭā)}]{\bn{বায়োডাটা} (bāyoḍāṭā) — a CV; a résumé}

\label{lesson:BN-C494-bayodata}



\begin{quote}
Before the new one: say the Bengali for to type, then the Bengali for unemployed.
\end{quote}

\subsection*{You'll want to know: \bn{বায়োডাটা}}
\textbf{\bn{বায়োডাটা}} — \emph{bāyoḍāṭā} — \textquotedblleft{}a CV; a résumé\textquotedblright{}.

\begin{quote}
\bn{আমি} \bn{চাকরির} \bn{জন্য} \bn{আমার} \bn{বায়োডাটা} \bn{পাঠালাম।}
\end{quote}

\begin{quote}
\emph{Āmi chākrir jônyô āmār bāyoḍāṭā pāṭhālām.}
\end{quote}

\begin{quote}
I sent my CV for the job.
\end{quote}

\begin{cousinweb}
\textbf{\bn{বায়োডাটা}} (\emph{bāyoḍāṭā}) is borrowed from Indian English \emph{bio-data}, a CV.
\end{cousinweb}

\begin{grammarlens}[title={What you've built}]
One more word for describing what you lived and felt.
\end{grammarlens}

\subsection*{Your turn}
\begin{itemize}
\raggedright
  \item \emph{Say it:} \emph{bāyoḍāṭā}
  \item \emph{Say it:} \emph{bāyoḍāṭā}, once more
  \item \emph{Say it:} \emph{Āmi chākrir jônyô āmār bāyoḍāṭā pāṭhālām.} — I sent my CV for the job.
  \item \emph{Recall:} say the Bengali for to type, then the Bengali for unemployed, then say \emph{bāyoḍāṭā} again
\end{itemize}

\begin{tcolorbox}[breakable,colback=teal!4,colframe=teal!35!black,title={Before you move on}]
What does \bn{বায়োডাটা} mean? (\textquotedblleft{}A CV; a résumé\textquotedblright{}.) Say it once more.
\end{tcolorbox}
\section[\texorpdfstring{\bn{শিক্ষানবিশ} (shikkhānôbish)}{শিক্ষানবিশ (shikkhānôbish)}]{\bn{শিক্ষানবিশ} (shikkhānôbish) — an apprentice; a trainee}

\label{lesson:BN-C494-shikkhanobish}



\begin{quote}
Before the new one: say the Bengali for unemployed, then the Bengali for a CV.
\end{quote}

\subsection*{You'll want to know: \bn{শিক্ষানবিশ}}
\textbf{\bn{শিক্ষানবিশ}} — \emph{shikkhānôbish} — \textquotedblleft{}an apprentice; a trainee\textquotedblright{}.

\begin{quote}
\bn{আমি} \bn{একটা} \bn{কারখানায়} \bn{শিক্ষানবিশ} \bn{হলাম।}
\end{quote}

\begin{quote}
\emph{Āmi ekṭā kārkhānāy shikkhānôbish hôlām.}
\end{quote}

\begin{quote}
I became an apprentice at a factory.
\end{quote}

\begin{cousinweb}
\textbf{\bn{শিক্ষানবিশ}} (\emph{shikkhānôbish}) joins Sanskrit \emph{śikṣā}, \textquotedblleft{}learning\textquotedblright{} (Monier-Williams), to an ending probably from Persian \emph{navīs}, \textquotedblleft{}one who writes\textquotedblright{}: a learner at a desk.
\end{cousinweb}

\begin{grammarlens}[title={What you've built}]
One more word for describing what you lived and felt.
\end{grammarlens}

\subsection*{Your turn}
\begin{itemize}
\raggedright
  \item \emph{Say it:} \emph{shikkhānôbish}
  \item \emph{Say it:} \emph{shikkhānôbish}, once more
  \item \emph{Say it:} \emph{Āmi ekṭā kārkhānāy shikkhānôbish hôlām.} — I became an apprentice at a factory.
  \item \emph{Recall:} say the Bengali for unemployed, then the Bengali for a CV, then say \emph{shikkhānôbish} again
\end{itemize}

\begin{tcolorbox}[breakable,colback=teal!4,colframe=teal!35!black,title={Before you move on}]
What does \bn{শিক্ষানবিশ} mean? (\textquotedblleft{}An apprentice; a trainee\textquotedblright{}.) Say it once more.
\end{tcolorbox}
\section[\texorpdfstring{\bn{ইন্টারভিউ} (inṭārbhiu)}{ইন্টারভিউ (inṭārbhiu)}]{\bn{ইন্টারভিউ} (inṭārbhiu) — a job interview}

\label{lesson:BN-C494-intarbhiu}



\begin{quote}
Before the new one: say the Bengali for a CV, then the Bengali for an apprentice.
\end{quote}

\subsection*{You'll want to know: \bn{ইন্টারভিউ}}
\textbf{\bn{ইন্টারভিউ}} — \emph{inṭārbhiu} — \textquotedblleft{}a job interview\textquotedblright{}.

\begin{quote}
\bn{কাল} \bn{আমার} \bn{একটা} \bn{ইন্টারভিউ} \bn{আছে।}
\end{quote}

\begin{quote}
\emph{Kāl āmār ekṭā inṭārbhiu āchhe.}
\end{quote}

\begin{quote}
I have an interview tomorrow.
\end{quote}

\begin{cousinweb}
\textbf{\bn{ইন্টারভিউ}} (\emph{inṭārbhiu}) is borrowed from English \emph{interview}.
\end{cousinweb}

\begin{grammarlens}[title={What you've built}]
One more word for describing what you lived and felt.
\end{grammarlens}

\subsection*{Your turn}
\begin{itemize}
\raggedright
  \item \emph{Say it:} \emph{inṭārbhiu}
  \item \emph{Say it:} \emph{inṭārbhiu}, once more
  \item \emph{Say it:} \emph{Kāl āmār ekṭā inṭārbhiu āchhe.} — I have an interview tomorrow.
  \item \emph{Recall:} say the Bengali for a CV, then the Bengali for an apprentice, then say \emph{inṭārbhiu} again
\end{itemize}

\begin{tcolorbox}[breakable,colback=teal!4,colframe=teal!35!black,title={Before you move on}]
What does \bn{ইন্টারভিউ} mean? (\textquotedblleft{}A job interview\textquotedblright{}.) Say it once more.
\end{tcolorbox}
\section[\texorpdfstring{\bn{যোগ} \bn{দেওয়া} (jog deoyā)}{যোগ দেওয়া (jog deoyā)}]{\bn{যোগ} \bn{দেওয়া} (jog deoyā) — to join a job; to join (a group)}

\label{lesson:BN-C494-jog-deoya}



\begin{quote}
Before the new one: say the Bengali for an apprentice, then the Bengali for a job interview.
\end{quote}

\subsection*{You'll want to know: \bn{যোগ} \bn{দেওয়া}}
\textbf{\bn{যোগ} \bn{দেওয়া}} — \emph{jog deoyā} — \textquotedblleft{}to join a job; to join (a group)\textquotedblright{}.

\begin{quote}
\bn{আমি} \bn{সোমবার} \bn{নতুন} \bn{অফিসে} \bn{যোগ} \bn{দিলাম।}
\end{quote}

\begin{quote}
\emph{Āmi sombār nôtun ôphise jog dilām.}
\end{quote}

\begin{quote}
I joined the new office on Monday.
\end{quote}

Past (he/she …): \bn{তিনি} \bn{নতুন} \bn{কোম্পানিতে} \bn{যোগ} \bn{দিলেন}, \emph{tini nôtun kompānite jog dilen}.

\begin{cousinweb}
\textbf{\bn{যোগ} \bn{দেওয়া}} (\emph{jog deoyā}) joins \bn{যোগ}, Sanskrit \emph{yoga}, \textquotedblleft{}joining\textquotedblright{} (Monier-Williams), to \bn{দেওয়া}, \textquotedblleft{}to give\textquotedblright{}: to give oneself to a place of work.
\end{cousinweb}

\begin{grammarlens}[title={What you've built}]
One more word for describing what you lived and felt.
\end{grammarlens}

\subsection*{Your turn}
\begin{itemize}
\raggedright
  \item \emph{Say it:} \emph{jog deoyā}
  \item \emph{Say it:} \emph{jog deoyā}, once more
  \item \emph{Say it:} \emph{Āmi sombār nôtun ôphise jog dilām.} — I joined the new office on Monday.
  \item \emph{Recall:} say the Bengali for an apprentice, then the Bengali for a job interview, then say \emph{jog deoyā} again
\end{itemize}

\begin{tcolorbox}[breakable,colback=teal!4,colframe=teal!35!black,title={Before you move on}]
What does \bn{যোগ} \bn{দেওয়া} mean? (\textquotedblleft{}To join a job; to join (a group)\textquotedblright{}.) Say it once more.
\end{tcolorbox}
